\addcontentsline{toc}{subsubsection}{Solution 35. The Subadditive--Supermultiplicative Squeeze}

\begin{dlsummary}[title={\bfseries Solution 35. The Subadditive--Supermultiplicative Squeeze},breakable]

\begin{proof}
The first inequality shows that $f$ is nondecreasing, and the second
shows that $f$ is subadditive. Since $f(x)+f(y)>0$, the last inequality
is equivalent to
\[
f(xy)\geq f(x)f(y). \tag{1}
\]
Thus $f$ is also supermultiplicative.

\medskip
\noindent
\textbf{Step 1. The anchor is a fixed point.}

Let $b=f(a)$. If $b>a$, monotonicity gives
\[
a=f(b)\geq f(a)=b>a,
\]
a contradiction. If $b<a$, then
\[
a=f(b)\leq f(a)=b<a,
\]
again a contradiction. Hence
\[
f(a)=a. \tag{2}
\]

\medskip
\noindent
\textbf{Step 2. Determine $f(1)$.}

Set $c=f(1)$. By (1), $c\geq c^2$, so $c\leq1$. For every positive
integer $n$, repeated use of (1) and (2) gives
\[
f(a^n)\geq f(a)^n=a^n.
\]
On the other hand, monotonicity and subadditivity give
\[
f(a^n)\leq f(\lceil a^n\rceil)
\leq \lceil a^n\rceil c.
\]
Therefore
\[
c\geq\frac{a^n}{\lceil a^n\rceil}\longrightarrow1,
\]
and so
\[
f(1)=1. \tag{3}
\]

\medskip
\noindent
\textbf{Step 3. Show that $f(x)\geq x$.}

Fix $x>0$. For all sufficiently large $n$, set
\[
m_n=\left\lfloor\frac{a^n}{x}\right\rfloor,
\qquad r_n=a^n-m_nx,
\]
so $m_n\geq1$ and $0\leq r_n<x$. By (1) and (2),
\[
a^n\leq f(a^n).
\]
If $r_n>0$, subadditivity, repeated subadditivity, and monotonicity give
\[
f(a^n)=f(m_nx+r_n)
\leq f(m_nx)+f(r_n)
\leq m_nf(x)+f(x).
\]
If $r_n=0$, the same final bound follows directly from
$f(m_nx)\leq m_nf(x)$. Thus
\[
f(x)\geq\frac{a^n}{m_n+1}
=\frac{a^n}{\left\lfloor a^n/x\right\rfloor+1}.
\]
Letting $n\to\infty$ yields
\[
f(x)\geq x. \tag{4}
\]

\medskip
\noindent
\textbf{Step 4. Force equality.}

Using (1), (3), and (4) applied to $1/x$, we obtain
\[
1=f(1)=f\!\left(x\cdot\frac1x\right)
\geq f(x)f\!\left(\frac1x\right)
\geq\frac{f(x)}x.
\]
Hence $f(x)\leq x$. Together with (4), this proves
\[
f(x)=x\qquad(x>0).
\]

The identity function satisfies the chain and $f(f(a))=a$ for every
$a>1$, so it is the unique solution.
\end{proof}

\end{dlsummary}